% ============================================================
% problem_analysis.tex - 问题分析
%
% 本章职责：建立对题目整体的认识，揭示各问之间的逻辑/数据/模型
% 依赖关系。每一问的**详细方法对比**放在 questionX.tex 的第一
% 小节，不要在本章展开。
% ============================================================

\section{问题分析}

\subsection{整体分析}

（用 3-5 句话说明题目的性质、核心难点、解决问题的整体思路。）

\subsection{问题间的依赖关系}

% 按建模阶段产出的依赖图，说明每一问之间是
% 依赖（Q(n+1) 用到 Q(n) 的模型或结果）/
% 扩展（Q(n+1) 在 Q(n) 基础上增加约束或变化条件）/
% 独立（完全独立求解）。
% 可用 tikz 绘制依赖关系图，或用如下的文字 + 表格方式说明。

\begin{table}[htbp]
  \centering
  \caption{问题间的依赖关系}
  \label{tab:dependency}
  \begin{tabular}{clll}
    \toprule
    \textbf{当前问题} & \textbf{依赖的前置问题} & \textbf{关系类型} & \textbf{传递内容} \\
    \midrule
    问题一 & —       & 基线       & —                         \\
    问题二 & 问题一  & 依赖       & 最优决策变量 $x^*$         \\
    问题三 & 问题二  & 扩展       & 增加鲁棒性约束             \\
    问题四 & —       & 独立       & —                         \\
    \bottomrule
  \end{tabular}
\end{table}

\subsection{各问分析概述}

对每一问给出问题类型判断（优化 / 预测 / 评价 / 分类 / 仿真）+ 关键
难点 + 思路方向（1 段即可，详细方法对比放在对应 questionX.tex）。

\textbf{问题一}：（问题类型 + 关键难点 + 思路方向）

\textbf{问题二}：（问题类型 + 关键难点 + 思路方向，如何依赖/扩展问题一）

% 按实际题目数量继续添加
